\section{The corner value: from exact rows to the first correction}
\label{sec:corner}

At a corner the absolute cardinal kernel becomes a positive Cauchy sum.
Its two meshes have sizes $n$ and $n+1$. The equal-mesh reference gives
the leading square; the mesh displacement supplies the constant defect
and the first correction.
Put $m=n+1$ and, for the rest of this section, write
\[
 \theta_j=j\pi/n,\quad \phi_k=k\pi/m,\qquad
 \varepsilon_j=\begin{cases}1/2,&j=0,n,\\1,&0<j<n,\end{cases}
 \quad \eta_k=\begin{cases}1/2,&k=0,m,\\1,&0<k<m.\end{cases}
\]
Every unspecified constant below is independent of $n\geq2$ and of the
row indices in their stated ranges. Set $A(z)=a(\log z+d)$, with the
constants in \eqref{eq:main-constants}.


The principal notation for this section is collected below. Row indices
and all implied constants have the ranges stated in the formulas.
\begin{center}
\begin{tabular}{@{}p{0.36\linewidth}p{0.58\linewidth}@{}}
\toprule
Symbol & Meaning\\
\midrule
$m=n+1$, $N=n+1/2$ & unequal mesh size and midpoint degree\\
$\theta_j$, $\phi_k$; $\varepsilon_j$, $\eta_k$ & angular nodes and endpoint half weights\\
$a,d,A(z)=a(\log z+d)$ & constants of \eqref{eq:main-constants} and logarithmic model\\
$\mathcal D_n$; $w(x),B(x)$ & diamond kernel; cotangent functions, with both poles removed from $B$\\
$x=r/m$, $s=m-r$ & position and complementary distance of row $r$\\
$T_{n,r},V_n,D_{m,r},\mathcal H_m$ & cotangent row sum, smooth sum, digamma row and its aggregate\\
$H_k=\sum_{j=1}^k1/j$ & harmonic numbers\\
\bottomrule
\end{tabular}
\end{center}
The product mesh size used in Section~\ref{sec:argmax} is instead
$\mathsf N=nm$.

\subsection{An exact row formula}
For $u\ne v$ let $f(u,v)=(1+\cos v)/|\cos u-\cos v|$. We first obtain
\begin{equation}\label{cu:positive}
 \Lambda_c(n)=\frac4{nm}
 \sum_{\substack{0\leq j,k\leq n\\j+k\text{ odd}}}
       \varepsilon_j\eta_k f(\theta_j,\phi_k)+\frac1m.
\end{equation}
To see this directly from \eqref{eq:cardinal-kernel}, define
$\mathcal{D}_n(u,v)=\sum_{|p|+|q|\leq n}\mathrm e^{\mathrm i(pu+qv)}$. Grouping positive and negative frequencies in each coordinate gives
\[
 \mathcal D_n(u,v)=\sum_{\substack{p,q\geq0\\p+q\leq n}}
       \nu_p\nu_q\cos(pu)\cos(qv)=G_n(u,v)+\cos(nu).
\]
For $B=(\cos\theta_j,\cos\phi_k)$ and $z=(-1,-1)$,
$T_p(-1)=(-1)^p$ therefore gives
$K_n(z,B)=\mathcal D_n(\pi-\theta_j,\pi-\phi_k)$.
Set $j'=n-j$ and $k'=m-k$. Then
$\pi-\theta_j=\theta_{j'}$, $\pi-\phi_k=\phi_{k'}$,
$(-1)^{n+j}=(-1)^{j'}$, and \eqref{eq:cardinal-kernel} becomes
\[
 L_{jk}(-1,-1)=w_{jk}
 \{\mathcal D_n(\theta_{j'},\phi_{k'})-(-1)^{j'}\}
 =w_{jk}G_n(\theta_{j'},\phi_{k'}).
\]
The reflection changes even checkerboard parity to odd and preserves
the endpoint weights. This is the link between the interpolation kernel
and the positive corner sum. With
$u_+=(u+v)/2$, $u_-=(v-u)/2$, finite geometric sums give
\[
 \mathcal{D}_n(u,v)=\frac{\sin((n+1)u_+)\sin((n+1)u_-)+\sin(nu_+)\sin(nu_-)}
                  {\sin u_+\sin u_-}.
\]
Indeed $(p,q)\mapsto(p+q,q-p)$ takes the diamond into the square
$[-n,n]^2$ with equal coordinate parity; summing its two parity classes
gives the displayed product. Quotients at zero denominators have their
continuous polynomial values. Reflection of node indices about either
nonnode corner turns the checkerboard odd and leaves its weights unchanged.
The corresponding Lebesgue sum is therefore
$\sum_{j+k\text{ odd}}2\varepsilon_j\eta_k
 |\mathcal{D}_n(\theta_j,\phi_k)-(-1)^j|/(nm)$.
For distinct mesh angles, the product-to-sum identities are
\[
 \sin(tu_+)\sin(tu_-)=\tfrac12\{\cos(tu)-\cos(tv)\},
 \qquad \sin u_+\sin u_-=\tfrac12(\cos u-\cos v).
\]
The four mesh cosine values are
\[
 \begin{gathered}
 \cos(n\theta_j)=(-1)^j,\quad
 \cos(m\theta_j)=(-1)^j\cos\theta_j,\\
 \cos(m\phi_k)=(-1)^k,\quad
 \cos(n\phi_k)=(-1)^k\cos\phi_k.
 \end{gathered}
\]
Since $j+k$ is odd, the numerator of $\mathcal D_n$ is
$\tfrac12(-1)^j(2+\cos\theta_j+\cos\phi_k)$.
Thus
\[
 \mathcal{D}_n(\theta_j,\phi_k)=(-1)^j
       \frac{2+\cos\theta_j+\cos\phi_k}{\cos\theta_j-\cos\phi_k},
 \qquad |\mathcal{D}_n(\theta_j,\phi_k)-(-1)^j|=2f(\theta_j,\phi_k).
\]
Coprimality of $n,m$ leaves only the collisions $(0,0)$ and $(n,m)$.
The odd checkerboard excludes the first. At the second,
$\mathcal{D}_n(\pi,\pi)=(-1)^n(2n+1)$, so the corrected absolute value $2n$
with weight $1/(2nm)$ contributes $1/m$. The other entries with $k=m$
vanish. This proves \eqref{cu:positive} and the equality of the two
nonnode corner values. The node corners have value one by cardinality.
For $n=2$, the nonzero weighted summands
$\varepsilon_j\eta_k f(\theta_j,\phi_k)$ at
$(j,k)=(0,1),(1,0),(1,2),(2,1)$ are respectively
$3/2,1,1,1/2$. Formula~\eqref{cu:positive} gives
$\Lambda_c(2)=(4/6)(3/2+1+1+1/2)+1/3=3$.


On row $k=r$ the checkerboard condition fixes $j\equiv1-r\pmod2$,
so $(-1)^j$ is constant along the row. The sign of the denominator
therefore determines its positive samples, making \eqref{cu:positive}
one dimensional. Introduce
\begin{equation}\label{cu:wB}
 w(x)=\cot(\pi x/2),\qquad
 B(x)=\cot(\pi x)-\frac1{\pi x}+\frac1{\pi(1-x)}.
\end{equation}
Both poles have been removed from $B$; it is $C^\infty$ across $[0,1]$,
with $B(0)=1/\pi$, $B(1)=-1/\pi$.
For example, with $\operatorname{sinc}z=\int_0^1\cos(tz)\,dt$,
\[
 \cot z-1/z=-z\,
 \frac{\int_0^1t^2\operatorname{sinc}(tz)\,dt}{\operatorname{sinc}z}.
\]
The denominator is nonzero near zero; differentiation of the finite
integral proves the smooth extension after subtracting the pole.
Reflection gives the extension at the other endpoint.
 For $1\leq r<m$ put
$x=r/m$, $s=m-r$, $\alpha_r=(1-x)/2$ and
\[
 T_{n,r}=\sum_{\ell=0}^{r-1}\cot\frac{\pi(\ell+\alpha_r)}n,
 \qquad V_n=\sum_{r=1}^{m-1}w(r/m)
                    \sum_{\ell=0}^{r-1}B((\ell+\alpha_r)/n).
\]
We define the digamma function $\psi$ by its convergent harmonic series:
\begin{equation}\label{cu:psi}
 \psi(z)=-\gamma+\sum_{k=0}^\infty
          \left(\frac1{k+1}-\frac1{k+z}\right),\qquad z>0,
\end{equation}
Telescoping this series gives
\[
 \psi(z+1)-\psi(z)=1/z,\qquad
 \psi(k+1)=H_k-\gamma,\qquad
 \psi(R+1/2)=\psi(1/2)+2H_{2R}-H_R,
\]
where $H_k=\sum_{j=1}^k1/j$ and $k,R$ are positive integers.
The bounds and half-integer values used below are proved in
Section~\ref{sec:asym}. In particular \eqref{cu:psi-est} gives
$H_k=\log k+\gamma+1/(2k)+O(k^{-2})$ by applying it at $k+1$.

Set
\begin{equation}\label{cu:D}
 D_{m,r}=\psi(r+\alpha_r)+\psi(s-\alpha_r)-\psi(m-\alpha_r)-\psi(\alpha_r),
 \qquad \mathcal H_m=\frac4{\pi m}\sum_{r=1}^{m-1}w(r/m)D_{m,r}.
\end{equation}
We claim the exact identities
\begin{equation}\label{cu:exact}
 \Lambda_c(n)=1-\frac{2n}{m}+\frac4{nm}\sum_{r=1}^{m-1}w(r/m)T_{n,r}
       =-1+\frac2m+\mathcal H_m+\frac{4V_n}{nm}.
\end{equation}
We now write row $k$ of \eqref{cu:positive} as $r$; its entries are
$|F_r(\theta_j)|$. Extend a row periodically
and sample its $n$ indices $j\equiv1-r\pmod2$ modulo $2n$. For
$v=r\pi/m$ its signed entries satisfy
\[
 F_r(u):=\frac{1+\cos v}{\cos u-\cos v}
   =\frac{w(r/m)}2\left(\cot\frac{v-u}{2}+\cot\frac{v+u}{2}\right),
 \qquad \sum_jF_r(j\pi/n)=n.
\]
The cotangent form uses
\[
 \begin{aligned}
 \cot A+\cot B&=\frac{\sin(A+B)}{\sin A\sin B},\\
 2\sin\frac{v-u}{2}\sin\frac{v+u}{2}&=\cos u-\cos v,\\
 \cot(v/2)\sin v&=1+\cos v.
 \end{aligned}
\]
The periodic row sum follows from
$\sum_{\ell=0}^{n-1}\cot(z+\ell\pi/n)=n\cot(nz)$, the logarithmic
derivative of the finite sine product; its right hand cotangent here is
$\tan(v/2)$. The positive samples have representatives
$j=-r+1,-r+3,\ldots,r-1$, since $r-1<nr/m<r$.
Their sum is $P_r=w(r/m)T_{n,r}$, by symmetry and the substitution
$j=r-1-2\ell$. Thus the absolute periodic row is $2P_r-n$.
Folding back to $0\leq j\leq n$ divides by two and produces exactly
$\varepsilon_j$. The folded interior row therefore contributes
$(4/(nm))\eta_r(2P_r-n)/2$ to the corner sum. For $r=0$ the periodic row is $n^2$, since the derivative
of the same cotangent identity gives
$\sum_{\ell=0}^{n-1}\csc^2((\ell+1/2)\pi/n)=n^2$.
Its weight $\eta_0=1/2$, together with the exceptional $1/m$, yields
$\Lambda_c=1/m+n/m+2\sum_{r=1}^{m-1}(2P_r-n)/(nm)$.
This is the first identity in \eqref{cu:exact}. Finally, removing both
poles in each cotangent and using $\psi(z+1)-\psi(z)=1/z$ gives
$T_{n,r}=nD_{m,r}/\pi+\sum_{\ell<r}B((\ell+\alpha_r)/n)$,
which proves the second identity.

\subsection{The square reference and its mesh defect}
The constant in the answer can be identified before calculating any
first corrections. Define
\[
 L_n=\frac2n\sum_{\substack{1\leq r\leq n\\r\text{ odd}}}
 d_r\csc\frac{r\pi}{2n},\qquad d_r=\begin{cases}1/2,&r=n,\\1,&r<n.\end{cases}
\]
The equal-mesh corner sum is exactly a square:
\begin{equation}\label{cu:square}
 \frac4{n^2}\sum_{\substack{0\leq j,k\leq n\\j+k\text{ odd}}}
 \varepsilon_j\varepsilon_k f(\theta_j,\theta_k)=L_n^2.
\end{equation}
Simultaneous reflection of $j,k$ replaces the numerator by $1-\cos\theta_k$,
so averaging replaces it by one. For $j>k$, put $s=j-k$, $r=j+k$ and
$c_r=\csc(r\pi/(2n))$. Both indices are odd and the denominator is
$2/(c_rc_s)$. The two orders cancel the factor two. For fixed $s<n$,
the values $r=s,2n-s$ have half weight and the intervening values weight
one; their sum after reflection $c_{2n-r}=c_r$ is
$c_s+2\sum_{s<r<n,\ r\text{ odd}}c_r+\boldsymbol1_{n\text{ odd}}c_n$.
The sole pair for $s=n$, when present, has weight $1/4$.
These are precisely the terms in $(\sum_{r\text{ odd}}d_rc_r)^2$,
proving \eqref{cu:square} including its endpoint.

Pairing the midpoints $y_r=(2r-1)\pi/(4n)$ gives, by
$\cot y+\tan y=2\csc(2y)$,
\[
 \begin{aligned}
 L_n&=\frac1n\sum_{r=1}^n\cot y_r\\
    &=a(2H_{2n}-H_n)+a\log(2/\pi)+O(n^{-2})=A(n)+O(n^{-2}),
 \end{aligned}
\]
Here $H_k=\sum_{j=1}^k j^{-1}$. For the second equality, $\cot y-1/y$
is smooth on $[0,\pi/2]$,
its integral is $\log(2/\pi)$, and the midpoint error is $O(n^{-2})$.
The last equality uses $H_k=\log k+\gamma+1/(2k)+O(k^{-2})$;
the reciprocal terms in $2H_{2n}-H_n$ cancel.

We next justify the limiting defect, including the endpoint rows. A
uniform comparison, with $\ell=|j-k|$ odd, is
\begin{equation}\label{cu:deform}
 |f(\theta_j,\phi_k)-f(\theta_j,\theta_k)|
                    \leq\frac{\pi^2n}{2\ell^2}.
\end{equation}
This bound follows from an exact difference. The cases $k=0$ or $j=n$ have zero
difference. Otherwise put $u=\theta_j$, $v_0=\theta_k$, $v_1=\phi_k$,
$\delta=v_0-v_1$, $V=(v_0+v_1)/2$, $S_i=(u+v_i)/2$ and
$Q_i=|u-v_i|/2$. The order of $u,v_i$ is unchanged. Exact subtraction gives
\[
 |f(u,v_1)-f(u,v_0)|=
 \frac{\cos^2(u/2)\sin V\sin(\delta/2)}
      {\sin S_0\sin S_1\sin Q_0\sin Q_1}.
\]
Since $\sin S_i\geq\cos(u/2)\sin(v_i/2)$, the quotient preceding
$\sin(\delta/2)$ and the $Q_i$ factors is at most
$\cot(v_0/2)+\cot(v_1/2)\leq\pi(m-k)/k$.
Here $\sin z\geq2z/\pi$ and $\cos z\leq\pi/2-z$ on $[0,\pi/2]$
suffice. Multiplying by $\sin(\delta/2)\leq\pi k/(2nm)$ gives
$\pi^2(m-k)/(2nm)$. If $k<j$, both $\sin Q_i\geq\ell/n$.
If $k>j$, put $h=n-j$; then
$\sin Q_0\geq\ell/n$, $\sin Q_1\geq[n(\ell-1)+h]/(nm)$ and
$m-k\leq h$. The resulting bound is
$\pi^2nh/[2\ell\{n(\ell-1)+h\}]\leq\pi^2n/(2\ell^2)$,
since $h\leq n$. This proves \eqref{cu:deform}.

Let $\Delta_n$ be $4/(nm)$ times the weighted sum of the differences in
\eqref{cu:deform}. The $k=n$ reference row is zero, so changing its weight
from $\eta_n$ to $\varepsilon_n$ has no effect. Equations
\eqref{cu:positive} and \eqref{cu:square} consequently give
\[
 \Lambda_c(n)-L_n^2=\Delta_n+(1-L_n^2)/m.
\]
The second term tends to zero. For fixed odd $\ell$ and $\sigma=\pm1$,
write $k=j+\sigma\ell$, $t=j/n$. On compact subintervals of $(0,1)$,
the cosine difference factors give uniformly
\[
 \frac{f(\theta_j,\phi_k)-f(\theta_j,\theta_k)}n
 \longrightarrow h_{\sigma,\ell}(t)
 =\frac{w(t)}\pi\left(\frac1{|\sigma\ell-t|}-\frac1\ell\right),
\]
because $\phi_k-\theta_j=\pi(\sigma\ell-t)/m$. Bound
\eqref{cu:deform} controls each normalized summand by $C/\ell^2$.
For fixed $\ell$, the rows with $t$ outside $[\delta,1-\delta]$ therefore
contribute at most $C(\delta+1/n)/\ell^2$; the same bound controls the
omitted limiting integrals. Interior Riemann sums, followed by
$\delta\downarrow0$, give the limit for that offset. The tails over
$\ell>M$ are uniformly bounded by $C\sum_{\ell>M}\ell^{-2}$, so they
may then be removed. Pairing orientations leaves
$1/(\ell-t)+1/(\ell+t)-2/\ell=2t^2/[\ell(\ell^2-t^2)]$.
Thus, on the full integer sequence,
\begin{equation}\label{cu:kappa}
 \Lambda_c(n)-A(n)^2\longrightarrow\kappa
\end{equation}
with exactly the integral \eqref{eq:main-kappa}. Its terms with $\ell\geq3$
converge uniformly up to both endpoints; for $\ell=1$ the zero of $w$ at
one cancels the pole. At zero $t^2w(t)\to0$. This also proves the claimed
continuous extension of the integrand.

\subsection{Uniform summation and the first correction}\label{sec:asym}
To obtain the first correction, we return to \eqref{cu:exact}, keeping
the harmonic and smooth parts together.
One elementary quadrature rule will control all remainders.

\begin{lemma}[Quadrature with endpoint logarithms]\label{cu:quadrature}
Suppose $g$ is continuous on $[0,1]$, twice continuously differentiable
inside, with $|g''(x)|\leq C(x^{-1}+(1-x)^{-1})$ there, and, for
$0<x\leq1/2$,
\[
 |g(x)-g(0)|+|g(1-x)-g(1)|\leq Cx(1+|\log x|).
\]
Then
\[
 \frac1m\sum_{r=1}^{m-1}g(r/m)=\int_0^1g(x)\,dx
       -\frac{g(0)+g(1)}{2m}+O(\log m/m^2).
\]
For $g\in C^2[0,1]$ the error is $O(m^{-2})$; its composite midpoint
rule has the same order.
\end{lemma}
\begin{proof}
On a cell $[u,u+h]$, two integrations by parts give
\[
 \frac h2\{g(u)+g(u+h)\}-\int_u^{u+h}g
 =\frac12\int_u^{u+h}(t-u)(u+h-t)g''(t)\,dt.
\]
For $h=1/m$ the two endpoint cells cost $O(h^2(1+|\log h|))$ by
the assumed moduli. On the others the kernel is at most $h^2/8$,
and $\int_h^{1-h}|g''|=O(\log m)$. For bounded second derivative,
the local error is at most $h^3\|g''\|_\infty/12$. Taylor's formula
at the midpoint gives the corresponding bound $h^3\|g''\|_\infty/24$.
\end{proof}

We also record the elementary estimates used below.
We use the following termwise-differentiation criterion: if a series of
continuously differentiable functions converges at one point of a compact
interval and the series of derivatives converges uniformly there, then
the sum is continuously differentiable and its derivative is that series.
Repeated application is valid when the successive derivative series have
the same property. For \eqref{cu:psi}, reciprocal powers give these uniform
bounds on every compact subinterval of $(0,\infty)$.
Differentiating
\eqref{cu:psi} on compact subsets of $z>0$ gives
\[
 \psi'(z)=\sum_{k\geq0}(z+k)^{-2},\quad
 \psi''(z)=-2\sum_{k\geq0}(z+k)^{-3},\quad
 |\psi''(z)|\leq z^{-2}+2z^{-3}.
\]
The unit-cell trapezoid identity applied to $(z+t)^{-1}$ and integral
comparison for $(z+t)^{-2}$ imply, uniformly for $z\geq1$,
\begin{equation}\label{cu:psi-est}
 \psi(z)=\log z-\frac1{2z}+O(z^{-2}),\qquad
 \psi'(z)=z^{-1}+O(z^{-2}).
\end{equation}
Indeed $\psi(z)=\lim_{K\to\infty}(\log K-\sum_{k<K}(z+k)^{-1})$,
and the trapezoid error is bounded by $\tfrac18\int_0^\infty
2(z+t)^{-3}dt$. This also proves the harmonic estimate used above.
The same series gives
$\psi(1/2)=-\gamma-2\log2$, $\psi'(1/2)=\pi^2/2$.
For the latter value, the classical paired partial fraction
$\pi\cot(\pi z)=z^{-1}+\sum_{k\geq1}((z-k)^{-1}+(z+k)^{-1})$
gives $\sum k^{-2}=\pi^2/6$ by its linear term at zero; removing even
indices gives $4\sum_{k\text{ odd}}k^{-2}=\pi^2/2$.
One short derivation of the partial fraction is the cosine series of
$\cos(zs)$ on $[0,\pi]$: its constant coefficient is
$2\sin(\pi z)/(\pi z)$ and its $k$th coefficient is
$2(-1)^kz\sin(\pi z)/(\pi(z^2-k^2))$. The coefficients are absolutely
summable; evaluating at $s=\pi$ gives the identity. All differentiated
series used here converge locally uniformly away from their poles.

Define the finite, endpoint-regular integrals
\[
 I_1=\int_0^1\left(\frac{w(x)}{1-x}-\frac ax\right)dx,
 \qquad I_2=\int_0^1xw(x)\,dx,\qquad I_B=\int_0^1xw(x)B(x)\,dx.
\]
The next proposition contains the only aggregate calculation needed.
\begin{proposition}\label{cu:aggregate}
There are constants $C_H,C_B$, independent of the degree, such that
\begin{align}
 \mathcal H_m&=A(m)^2+C_H+
       \frac{-a^2\log m+\beta_H}{m}+O(\log m/m^2),\label{cu:H}\\
 \frac{4V_n}{nm}&=C_B+\frac{\beta_B}{n}+O(n^{-2}),\label{cu:V}\\
 \beta_H&=-2-a^2(\gamma+2\log2)-a(I_1+I_2),\notag\\
 C_B&=4\int_0^1w(t)\int_0^tB(x)\,dx\,dt,
 \qquad \beta_B=-a^2+2I_B+aI_2.\notag
\end{align}
\end{proposition}
\begin{proof}
First center the harmonic row at its exact integer and half integer
arguments, before expanding. Taylor's formula gives, uniformly for
$r+s=m$, $x=r/m$, $\alpha=(1-x)/2$,
\begin{equation}\label{cu:center}
 \begin{split}
 D_{m,r}={}&\psi(r+1/2)+\psi(s)-\psi(\alpha)-\psi(m)\\
 &-\frac{r\psi'(r+1/2)+s\psi'(s)-s\psi'(m)}{2m}+O(m^{-2}).
 \end{split}
\end{equation}
In fact the three shifts are respectively $r/(2m),s/(2m),s/(2m)$.
For a shift $h$ at $z$ the error is at most
$h^2\{(z-h)^{-2}+2(z-h)^{-3}\}/2$.
The lower arguments are at least $r$, $3s/4$, $3m/4$, respectively;
hence each error is $O(m^{-2})$, even for $r=1$ or $s=1$.
Since $w(r/m)\leq am/r$, its weighted normalized error is
$O(\log m/m^2)$.

Two summable endpoint sequences must still be retained. Put
$e_r=\psi(r+1/2)-\log r$, $u_r=r\psi'(r+1/2)-1$.
The midpoint rule for $t^{-2}$ on $[r,\infty)$ gives
$\psi'(r+1/2)=1/r+O(r^{-3})$: the local errors sum as
$O(\sum_{k\geq0}(r+k)^{-4})$. Integration from infinity, using
\eqref{cu:psi-est}, then gives $e_r=O(r^{-2})$ as well as
$u_r=O(r^{-2})$. Thus $S_e=\sum e_r/r$ and $S_u=\sum u_r/r$ converge.
Finite recurrence summation gives
\[
 \sum_{r=1}^R\psi'(r+1/2)
 =(R+1/2)\psi'(R+1/2)+\psi(R+1/2)
               -\psi(1/2)-\tfrac12\psi'(1/2).
\]
This follows by summing
$\psi'(r+1/2)=\psi'(1/2)-\sum_{k<r}(k+1/2)^{-2}$ and writing
$k=(k+1/2)-1/2$ in the finite triangle. Subtracting $H_R$ and passing
to the limit gives $S_u=1+2\log2-\pi^2/4$.

Let $d_0=\gamma+2\log2$, $q_0=\pi^2/4-1$ and
$F(x)=\log x+\log(1-x)-\psi((1-x)/2)$. Substitution of
\eqref{cu:psi-est} into \eqref{cu:center}, keeping $e_r,u_r$ intact, yields
\begin{equation}\label{cu:row-expand}
 D_{m,r}=\log m+F(x)+e_r-\frac1{2m(1-x)}-\frac{x+u_r}{2m}+Q_{m,r},
 \quad |Q_{m,r}|\leq C\{s^{-2}+(ms)^{-1}+m^{-2}\}.
\end{equation}
Indeed $\psi(s)-\psi(m)$ contributes
$\log(s/m)-1/(2s)+1/(2m)$, whereas the derivative bracket is
$1+x+u_r+O(1/s+1/m)$; its constant cancels $1/(2m)$.
To sum the remainder, use $w(r/m)\leq am/r$ on the first half and
$w(r/m)\leq Cs/m$ on $r\geq m/2$. They give respectively bounds
$C/(m^2r)$ and $C/(m^2s)$ after normalization by $1/m$.
Consequently $m^{-1}\sum w|Q|=O(\log m/m^2)$.

Write $v(x)=w(x)-a/x$, smoothly extended by $v(0)=0$; $|v(x)|\leq Cx$.
The sequence bounds and their $O(m^{-2})$ tails show
$m^{-1}\sum_{r<m}w(r/m)e_r=aS_e+O(\log m/m^2)$,
and the same assertion holds with $u_r,S_u$. Each $v$ contribution is
bounded by $Cm^{-2}\sum_{r<m}r^{-1}$. Ordinary trapezoid estimates for
the smooth functions $w/(1-x)-a/x$ and $xw$ also give
\[
 \frac1m\sum_{r<m}\frac{w(r/m)}{1-r/m}
     =a\log m+a\gamma+I_1+O(m^{-1}),\qquad
 \frac1m\sum_{r<m}(r/m)w(r/m)=I_2+O(m^{-1}).
\]
Thus, for $M_m=m^{-1}\sum_{r<m}w(r/m)\{\log m+F(r/m)\}$,
\begin{equation}\label{cu:H-to-M}
 \mathcal H_m=2a\left\{M_m+aS_e-
    \frac{a\log m+a\gamma+I_1+I_2+aS_u}{2m}\right\}
                      +O(\log m/m^2).
\end{equation}

It remains to compute the endpoint correction of $M_m$. Subtract its
explicit singularity:
$R(x)=w(x)F(x)-a(\log x+d_0)/x$.
Near zero $J(x)=\log(1-x)-\psi((1-x)/2)$ is analytic, with
$J(0)=d_0$, $J'(0)=q_0$, and
\[
 R(x)=a\frac{J(x)-d_0}{x}+v(x)\{\log x+J(x)\}.
\]
Near one, write $y=1-x$ and use
$\psi(y/2)=-2/y+\psi(1+y/2)$, $w(1-y)=\tan(\pi y/2)$.
These analytic factors prove
$R(0)=aq_0$, $R(1)=\pi-ad_0$ and exactly the derivative and endpoint
modulus bounds of Lemma~\ref{cu:quadrature}. Its conclusion applies.
Also $v(1)=-a$, $\int_0^1v=a\log(2/\pi)$, so smooth quadrature gives
\begin{equation}\label{cu:w-sum}
 \frac1m\sum_{r<m}w(r/m)=a\{\log m+\gamma+\log(2/\pi)\}+O(m^{-2}).
\end{equation}
The $-a/(2m)$ from $aH_{m-1}$ cancels the endpoint correction for $v$.
Finally, unit-cell trapezoids for $\log t/t$ give a real constant $C_K$ with
\[
 H_{m-1}=\log m+\gamma-\frac1{2m}+O(m^{-2}),\qquad
 \sum_{r<m}\frac{\log r}{r}=\tfrac12\log^2m+C_K
                   -\frac{\log m}{2m}+O(\log m/m^2).
\]
For the second formula the local errors have a convergent sum because
$(\log t/t)''=(2\log t-3)/t^3$; the tail is $O(\log m/m^2)$.
Express $M_m$ as $\log m$ times \eqref{cu:w-sum}, plus
$a\sum_{r<m}(\log r)/r+a(d_0-\log m)H_{m-1}$, plus $m^{-1}\sum R$.
The logarithmic reciprocal terms cancel, leaving
\[
 M_m=\tfrac a2\log^2m+ad\log m+C_M
          -\frac{\pi+aq_0}{2m}+O(\log m/m^2)
\]
for a constant $C_M$. Inserting this into \eqref{cu:H-to-M} and using
$q_0+S_u=2\log2$ proves \eqref{cu:H} and its stated $\beta_H$;
all degree-independent terms are absorbed into $C_H$.

For the smooth part, the geometry of the cells supplies the correction.
The points $(\ell+\alpha)/n$, $0\leq\ell<r$, are the midpoints of
cells whose union is exactly $[-t/(2n),t+t/(2n)]$, where $t=r/m$.
On $[-1/4,5/4]$ the function $B$ and its first two derivatives are bounded.
The midpoint error in the sum is $O(r/n^2)=O(t/n)$; integrating the two
short extensions of $[0,t]$ therefore gives uniformly
\[
 \sum_{\ell<r}B((\ell+\alpha)/n)
   =nI(t)+\tfrac t2\{B(t)+B(0)\}+O(t/n),\qquad I(t)=\int_0^tB(x)\,dx.
\]
Put $W(t)=tw(t)$, smoothly extended by $W(0)=a$.
The functions $F_B(t)=w(t)I(t)=W(t)\int_0^1B(st)\,ds$ and
$G_B(t)=W(t)\{B(t)+1/\pi\}$ are smooth on $[0,1]$.
Their endpoint values are $F_B(0)=a^2/2$, $G_B(0)=a^2$,
$F_B(1)=G_B(1)=0$. Weighting the row expansion, using boundedness
of $W$, and applying smooth quadrature yields
\[
 \begin{aligned}
 \frac{4V_n}{nm}
 &=\frac4m\sum_{r<m}F_B(r/m)+\frac2{nm}\sum_{r<m}G_B(r/m)+O(n^{-2})\\
 &=4\int_0^1F_B-\frac{a^2}{m}
       +\frac2n\int_0^1G_B-\frac{a^2}{nm}+O(n^{-2}).
 \end{aligned}
\]
Since $1/m+1/(nm)=1/n$ and $\int G_B=I_B+I_2/\pi$, this is
\eqref{cu:V}, with its stated constants. The calculation keeps the last
row as well as both outer endpoint corrections.
\end{proof}

\begin{proof}[Proof of Theorem~\ref{thm:main-expansion}]
Only one integral identity is needed to combine the two corrections:
\begin{equation}\label{cu:cancel}
 2I_B-aI_1=a^2-1+a^2\log(\pi/2).
\end{equation}
The double-angle identity $\cot(\pi x)=(w-w^{-1})/2$ makes its
combined integrand $x(w^2-1)-2aw+a^2/x$.
Since $w'=-(1+w^2)/a$, this equals
$-a(xw)'-a(w-a/x)-2x$. Integration gives \eqref{cu:cancel}, using
$\lim_{x\downarrow0}xw(x)=a$, $w(1)=0$ and
$\int_0^1(w-a/x)=a\log(2/\pi)$.

Replacing $\beta_B/n$ by $\beta_B/m$ costs $O(m^{-2})$.
The scalar $2/m$ in \eqref{cu:exact} now combines with the two
corrections as
\[
 2+\beta_H+\beta_B
 =-a^2-a^2(\gamma+2\log2)+2I_B-aI_1=-1-a^2d.
\]
Consequently, for a constant $C=-1+C_H+C_B$,
\[
 \Lambda_c(n)=A(m)^2+C-
       \frac{a^2\log m+1+a^2d}{m}+O(\log m/m^2).
\]
The limit \eqref{cu:kappa}, and $A(m)^2-A(n)^2\to0$, identify $C=\kappa$.
For $N=m-1/2$, Taylor's formula gives uniformly
\[
 A(N)^2=A(m)^2-\frac{a^2\log m+a^2d}{m}+O(\log m/m^2),
 \qquad N^{-1}=m^{-1}+O(m^{-2}).
\]
This proves $\Lambda_c(n)=A(N)^2+\kappa-N^{-1}+O(\log N/N^2)$.
Expanding the square gives \eqref{eq:main-expansion} with precisely the
constants \eqref{eq:main-constants}. Every bound holds for all $n\geq2$,
so the remainder is uniform across both parities. The corner proof is
entirely analytic; its application to $\Lambda_n$ uses
Theorem~\ref{thm:main-location}.
\end{proof}
